% ==============================================================================
% FFGFT: Dokument 245 (DE)
% Titel: FFGFT in der RA 2.1 Architektur — Diagnose-Eintrag
% Autor: Johann Pascher · Mai 2026
% RA-Nomenklatur: Guevara Calderón, J.T. (2026). RA 2.1. IPI-Arbeitsdokument.
% ==============================================================================

\subsection*{Zeichenerklärung}

{\small
\begin{tabular}{p{3.2cm}p{10.5cm}}
\toprule
\textbf{Symbol} & \textbf{Bedeutung} \\
\midrule
$\xi = 4/30000$ & FFGFT dimensionsloser Wicklungsparameter; geometrisch fixiert über $\kappa=7$ (nicht gefittet) \\
$T^4$ & 4-dimensionaler Torus $(\mathbb{R}/2\pi\mathbb{Z})^4$; ontologisches Primitiv von FFGFT \\
$n_\theta, n_\phi, n_3, n_4$ & Ganzzahlige Wicklungszahlen in den vier Torusrichtungen \\
$\theta_4$ & Kontinuierliche vierte Phasenkoordinate auf $T^4$ \\
$\varepsilon \approx 0.001$ & Nicht-Abschluss von $T^4$; generative strukturelle Residue \\
$\tilde{T}\cdot m = 1$ & Grundlegende T0-Relation (Zeit-Masse-Dualität) \\
$D_f = 3 - \xi$ & Fraktale Dimension des physikalischen Raums ($\approx 2{,}99987$) \\
$\kappa = 7$ & Eindeutiger Wicklungsexponent aus drei gemessenen Massenverhältnissen \\
$L^2(T^4)$ & FFGFT Hilbertraum; quadratintegrierbare Funktionen auf $T^4$ \\
$\alpha$ & Feinstrukturkonstante; aus $\xi$ hergeleitet (Dok.~009/137) \\
$\Lambda_\text{FFGFT}$ & Kosmologische Konstante: $1{,}34\times10^{-122}$ Planck-Einheiten (Dok.~182) \\
$H_0$ & Hubble-Konstante aus $\xi$: $66{,}2\,\text{km/s/Mpc}$, $-1{,}9\%$ von Planck (Dok.~182) \\
$\tau\cdot\mu$ & Kohärenzzeit--Informationsmasse-Produkt $= \hbar/c^2$; exakt \\
RA4--RA1 & Schichten der Darstellungsarchitektur, eingeführt von J.T.~Guevara Calderón (2026), \textit{RA 2.1}, IPI-Arbeitsdokument, Mai 2026 \\
\bottomrule
\end{tabular}
}

\vspace{0.6em}\hrule\vspace{0.5em}

\begin{center}
{\large\bfseries FFGFT Diagnose-Eintrag — RA 2.1 Format}\\[0.2em]
{\normalsize Fundamentale Fraktale Geometrische Feldtheorie (T0 / Zeit-Masse-Dualität)}\\[0.1em]
{\small Johann Pascher · ORCID 0009-0000-6518-4064 · Mai 2026}\\[0.1em]
{\footnotesize RA-Nomenklatur: Guevara Calderón, J.T. (2026). \textit{Representational Architecture 2.1}. IPI-Arbeitsdokument, Mai 2026.}
\end{center}

\vspace{0.3em}\hrule\vspace{0.3em}

\renewcommand{\arraystretch}{1.4}
{\small
\begin{longtable}{p{2.0cm}p{4.4cm}p{4.4cm}p{2.55cm}}

\toprule
\textbf{Schicht} & \textbf{FFGFT Inhalt} & \textbf{Standardmodell / $\Lambda$CDM} & \textbf{RA-Status} \\
\midrule
\endfirsthead
\toprule
\textbf{Schicht} & \textbf{FFGFT Inhalt} & \textbf{Standardmodell / $\Lambda$CDM} & \textbf{RA-Status} \\
\midrule
\endhead
\midrule\multicolumn{4}{r}{\footnotesize\itshape Fortsetzung\ldots}\\\endfoot
\bottomrule\endlastfoot
\sloppy\raggedright

\textbf{RA4}\\Ontol. &
4D-Torus $T^4$ mit Wicklungsstruktur $(n_\theta, n_\phi, n_3, n_4)$, kontinuierlicher Phase $\theta_4$ und Nicht-Abschluss $\varepsilon \approx 0{,}001$. Einziger Parameter $\xi = 4/30000$ geometrisch fixiert aus drei Massenverhältnissen über $\kappa = 7$. Grundrelation $\tilde{T}\cdot m = 1$. &
Quantenfelder auf Minkowski-Raumzeit $\mathbb{R}^{3,1}$. Mehrere unabhängige Kopplungskonstanten, Massen und Mischungswinkel als Eingaben. $\Lambda$ mit QFT-Vakuumenergie identifiziert (Zeldovich Schritt~3). &
\textbf{RA4 Divergenz.} Primitiv ersetzt. Zeldovich Schritt~3 nicht gemacht; kein Fein-\allowbreak abstimmungs-\allowbreak problem. \\

\midrule

\textbf{RA3}\\Zuläss. &
Nur mit $\xi$ konsistente Wicklungskonfigurationen physikalisch zulässig. Fraktale Dimension $D_f = 3-\xi \approx 2{,}99987$ definiert zugängliche topologische Sektoren. Z3$^3$-Bedingung beschränkt auf genau drei Teilchengenerationen. Keine weiteren Axiome nötig. &
Alle renormierbaren Feldkonfigurationen zulässig. Teilcheninhalt und Generationsanzahl sind experimentelle Eingaben, keine abgeleiteten Bedingungen. &
\textbf{RA3 Einschränkung.} $\xi$ schränkt Zulässigkeit maximal ein. SM lässt sie offen; Generationen ungeklärt. \\

\midrule

\textbf{RA2}\\Darst. &
$L^2(T^4)$ als Zustandsraum. Fourier-Moden mit Wicklungszahlen $\mathbf{n}\in\mathbb{Z}^4$ als natürliche Basis. QM-Operatoralgebra und Born-Regel entstehen als Projektion aus der Torus-Geometrie (Dok.~230/231, Hilbertraum-Brücke). Kein separates Quantisierungspostulat. &
Fock-Raum über Minkowski. Operatoralgebra und Born-Regel werden unabhängig von der Raumzeitstruktur postuliert. Quantisierung ist eine separate Eingabe. &
\textbf{RA2 Herleitung.} QM-Formalismus aus $T^4$-Geometrie hergeleitet. Im SM wird er postuliert. \\

\midrule

\textbf{RA1.5}\\Operat. &
Alles aus $\xi$ allein hergeleitet, \textbf{kein Fitting}:
{\small\begin{itemize}[leftmargin=*,nosep,itemsep=1pt]
  \item Lepton- und Quarkmassen via $\kappa=7$ Rekursion; Übereinstimmung $<1\%$
  \item Feinstrukturkonstante $\alpha$ (Dok.~009/137)
  \item CMB-Temperatur; Abweichung $\Delta = 0{,}0002\%$ (Dok.~025)
  \item $\Lambda_\text{FFGFT} = 1{,}34\times10^{-122}$ Planck-Einheiten (Dok.~182)
  \item $H_0 = 66{,}2\,\text{km/s/Mpc}$ ($-1{,}9\%$ von Planck)
  \item $\tau\cdot\mu = \hbar/c^2$ (exakt)
\end{itemize}} &
Alle genannten Größen sind unabhängige experimentelle Eingaben. $\Lambda$ erfordert Feinabstimmung über 120 Größenordnungen gegenüber der QFT-Nullpunktsenergie. &
\textbf{RA1.5 Unifikation.} Alle SM-Konstanten aus einem Parameter. SM behandelt sie als unabhängige Eingaben. \\

\midrule

\textbf{RA1}\\Ausdr. &
Konkrete numerische Werte für alle SM-Konstanten und kosmologischen Parameter, hergeleitet nicht gefittet. Veröffentlicht in 149+ Dokumenten; GitHub: jpascher/T0-Time-Mass-Duality; Zenodo v1.1.0, DOI 10.5281/zenodo.20117635. &
Parameterwerte experimentell gemessen und manuell eingesetzt. Keine vereinheitlichende Herleitung. SM hat $\approx\!20$ freie Parameter; $\Lambda$CDM weitere. &
\textbf{RA1 Kompression.} 1 Parameter (FFGFT) vs.\ $\sim$20 (SM). \\

\end{longtable}
}

\vspace{0.3em}\hrule\vspace{0.4em}

\noindent\textbf{Falsifizierungskriterien (RA1-Ebene):}
\begin{itemize}[nosep,itemsep=1pt]
  \item Jedes Lepton- oder Quark-Massenverhältnis das mit der $\xi$-Leiter bei $\kappa=7$ inkonsistent ist (Toleranz $<1\%$).
  \item CMB-Abweichung die den Korrekturterm $(275/4)\cdot\xi \approx 9{,}2\times10^{-3}$ überschreitet (Dok.~025).
  \item Ein Wert der kosmologischen Konstante außerhalb des in Dok.~182 abgeleiteten Bereichs.
  \item Jede Wicklungsquantenzahl die keiner ganzen Zahl in $\mathbb{Z}^4$ zugeordnet werden kann.
\end{itemize}

\vspace{0.3em}
\noindent\textbf{Schlüsseldokumente:}
Dok.~009 ($\xi$, $\kappa=7$) ·
Dok.~013 (SI-Einheiten) ·
Dok.~025 (CMB) ·
Dok.~182 ($\Lambda$, $H_0$) ·
Dok.~230/231 (Hilbertraum) ·
\href{https://github.com/jpascher/T0-Time-Mass-Duality}{GitHub} ·
\href{https://doi.org/10.5281/zenodo.20117635}{Zenodo v1.1.0}
